\section{Introduction}
\label{sec:intro}
Recognising an action from video means combining what each frame shows
with how it changes over time. We study this on a ten-class workout subset
of UCF101~\cite{soomro2012ucf101}, where classes such as \emph{PushUps} and
\emph{WallPushups}, or \emph{JumpingJack} and \emph{JumpRope}, share poses
and scenes but differ in motion. We ask three questions.
(i)~Where should a CNN combine the frames: after classifying each one,
after pooling their features, at the input, or throughout a 3D network
(Task~4.1)?
(ii)~How much do the original splits, in which the same people appear in
training and evaluation, inflate these accuracies (Task~4.2.1)?
(iii)~Does an optical-flow stream help once identity cues are removed
(Task~4.2.2)?
So that the comparison is about the fusion point rather than the network,
all models share one backbone and one training recipe.

\paragraph{Related work.}
Karpathy \etal~\cite{karpathy2014largescale} compared single-frame, early,
late, and slow fusion of 2D CNNs and found single-frame models
surprisingly strong. C3D~\cite{tran2015c3d} learns spatiotemporal features
with $3{\times}3{\times}3$ kernels, and I3D~\cite{carreira2017i3d}
initialises such kernels by inflating trained 2D filters. Two-stream
networks~\cite{simonyan2014twostream} add a temporal stream on stacked
optical flow and average the two streams' softmax scores; TSN~\cite{wang2016tsn}
initialises the flow network from RGB weights. We apply these ideas at small
scale, training on only 500 videos.
